\documentclass[11pt]{article}
\usepackage[margin=2.6cm]{geometry}
\usepackage{amsmath,amssymb,amsthm}
\usepackage{enumitem}
\usepackage{xspace}
\usepackage[T1]{fontenc}
\setlist{itemsep=1pt,topsep=3pt}

% the name of the system: change it here
\newcommand{\sys}{\textsf{Unit}\xspace}

\newcommand{\kw}[1]{\texttt{#1}}
\newcommand{\U}{\mathsf{U}}
\newcommand{\Id}{\mathsf{Id}}

\theoremstyle{remark}
\newtheorem{remark}{Remark}

\title{\sys: a small dependent type checker\\with segments as modules}
\author{}
\date{17 October 2026}

\begin{document}
\maketitle

\begin{abstract}
\noindent
\sys is a prototype proof checker for a dependent type theory with
$\Pi$, $\Sigma$, records, labelled sums, an identity type with its
eliminator, and a hierarchy of cumulative universes indexed by level
expressions. Its module system is based on the segments of
\textsc{Automath}, in the form given in~\cite{Wang}: a module is a list of
declarations, its parameters are the declarations without a body, and
an instance only binds the parameters, without copying the body. The same
segments also serve as record types. The prototype has been used to
formalise a model of univalent type theory over posets and simplicial
sets, about 200 files. This note describes the type theory, the module
system, and the choices made to keep checking predictable and fast. The
name \sys refers to the units, the declarations of which every file,
segment and local block is a list.
\end{abstract}

\tableofcontents

%======================================================================
\section{Overview}

A file is a list of \emph{units}. A unit is a declaration of one of the
following kinds:
\begin{center}
\begin{tabular}{ll}
\kw{x (tel) : A ;} & a parameter (at the top level: a primitive notion)\\
\kw{x (tel) : A = e ;} & a definition\\
\kw{data D (tel) = c$_1$ (tel$_1$) | \dots ;} & a labelled sum (a definition)\\
\kw{level a b ;} & level variables\\
\kw{constraint a < b ;} & constraints between levels\\
\kw{module M (tel) \{ units \} ;} & a segment\\
\kw{module N = M[args] ;} & an instance of a segment\\
\kw{open M ;}\quad \kw{import F ;} & scoping\\
\kw{opaque f ;}\quad\kw{transparent f ;} & the unfolding strategy
\end{tabular}
\end{center}
Checking a list of units threads a \emph{context of units}: each unit is
checked in the context of the units before it, and is then added to it.
A definition stays a definition in the context: its body is available for
computation. The bodies of \kw{let} and \kw{where} are lists of units as
well. Consecutive definitions form one mutually recursive group.

Two design principles run through the system. First, \emph{nothing is
guessed}: there is no unification, and the few arguments that may be
omitted are found by a deterministic matching procedure
(Section~\ref{sec:hidden}). Second, \emph{nothing is copied}: a
definition has one identity wherever it is used, and two uses are
compared by that identity before anything is unfolded
(Section~\ref{sec:conv}).

%======================================================================
\section{The type theory}\label{sec:tt}

\subsection{Terms}

The terms are those of a dependent type theory with definitions:
\begin{center}
\begin{tabular}{ll}
\kw{U l} & the universe at level \kw{l}\\
\kw{(x : A) -> B},\quad \kw{\textbackslash\ x -> t},\quad \kw{t u} & dependent functions\\
\kw{(x : A) * B},\quad \kw{(t, u)},\quad \kw{t.1},\quad \kw{t.2} & dependent pairs\\
\kw{sig \{ l$_1$ : A$_1$ ; \dots \}},\quad \kw{struct \{ l$_1$ = t$_1$ ; \dots \}},\quad \kw{t.l} & records\\
\kw{c t$_1$ \dots t$_n$},\quad \kw{split \{ c$_1$ x$_1$\dots -> t$_1$ ; \dots \}} & constructors and case analysis\\
\kw{Id A a b},\quad \kw{refl a},\quad \kw{idJ A a C d b p} & the identity type\\
\kw{[a] A},\quad \kw{<a> t},\quad \kw{t @l} & products over levels\\
\kw{[a < b] A} & products over constraints\\
\kw{let \{ units \} in t},\quad \kw{t where \{ units \}} & local definitions\\
\kw{M[args].x} & a member of a segment instance
\end{tabular}
\end{center}

\paragraph{Labelled sums.} A \kw{data} declaration introduces a type with
named constructors, each with a telescope of arguments; the declaration
may be recursive. A function out of a sum is written \kw{split}, a list
of branches, one per constructor; \kw{split} is not a term on its own but
is checked against a function type whose domain is a sum, as in
\begin{verbatim}
  data N = zero | suc (n : N) ;
  plus (n : N) : N -> N = split { zero -> n ; suc m -> suc (plus n m) } ;
\end{verbatim}
Recursion is through the name of the definition. There are no indexed
inductive families. Neither termination nor positivity is checked.

\paragraph{The identity type.} \kw{Id A a b} is a primitive type with
\kw{refl a : Id A a a} and the eliminator
\[
\frac{C : (y:A)\to\Id\,A\,a\,y\to\U\,l\qquad d : C\,a\,(\mathsf{refl}\,a)\qquad p:\Id\,A\,a\,b}
     {\mathsf{idJ}\,A\,a\,C\,d\,b\,p : C\,b\,p}
\]
with the computation rule $\mathsf{idJ}\,A\,a\,C\,d\,a\,(\mathsf{refl}\,a)=d$.
Since there is no pattern matching on \kw{refl}, there is no way to derive
uniqueness of identity proofs: the theory is compatible with univalence,
which a development may assume as a parameter
(Section~\ref{sec:params}).

\subsection{Universes and levels}\label{sec:univ}

The universes follow the system of~\cite{BCDE}.
Levels are expressions
\[
l \;::=\; a \;\mid\; l^{+} \;\mid\; l\vee m
\]
in level variables $a$, a successor and a maximum. There is no level $0$ for
the user. Level variables are declared by \kw{level a}, and may be
constrained by \kw{constraint} units ($l<m$, $l\le m$, $l=m$).

The universes are \emph{\`a la Russell}: an element of \kw{U l} is itself a
type. The rules are
\begin{itemize}
\item $\U\,l : \U\,l^{+}$;
\item $\Pi$, $\Sigma$, records and $\Id$ lie in the maximum of the levels
  of their components;
\item cumulativity: if $A:\U\,l$ and $l\le m$ then $A:\U\,m$.
  Cumulativity applies to the whole type and not under a binder: a
  function into $\U\,l$ is not silently a function into $\U\,m$, and one
  writes $\lambda b.\,T\,b$ when needed;
\item a sum declared without a universe, like \kw{N} above, is
  \emph{small}: it lies in every universe. A sum may also be declared at a
  given level, \kw{data D : U l = \dots}.
\end{itemize}

\paragraph{Products over levels and constraints.} A type may be
polymorphic in a level: \kw{[a] A} is the type of the \kw{<a> t}, applied
to a level by \kw{t @l}. A type may also depend on a constraint:
\kw{[a < b] A}. A constraint product is introduced and eliminated
silently: a term of type $[\psi]A$ is used as a term of type $A$ wherever
$\psi$ holds in the current context. Both kinds of products are
\emph{large}: they lie in no universe.

\paragraph{Deciding constraints.} Whether a constraint follows from those
of the context is decided by the algorithm of Bezem and Coquand for
join-semilattices with an inflationary endomorphism: constraints become
Horn clauses, and the question is answered by computing a least model by
forward reasoning~\cite{BC}. The same algorithm detects \emph{loops}, contexts
whose constraints force $l<l$. A context with a loop is inconsistent, and
in it every two terms are convertible. In practice this signals an error
in the constraints, and the checker warns.

\subsection{Definitional equality}\label{sec:conv}

Definitional equality is generated by
\begin{itemize}
\item $\beta$ for functions, pairs, records, level abstractions and
  \kw{split} applied to a constructor;
\item the computation rule of \kw{idJ} on \kw{refl};
\item the unfolding of definitions (but see \kw{opaque} below);
\item $\eta$ for functions, pairs, records and elements of segment types.
\end{itemize}

The checker decides it as follows. Two values are first compared
\emph{structurally}, without any computation. A use of a definition is
represented by its name together with the values of the variables its
body depends on (its \emph{arguments}, computed once, when the definition
is declared); two uses of the same definition with equal arguments are
equal without unfolding. Only when this fails are both sides brought to
head normal form, unfolding definitions, and compared component by
component.

\paragraph{Records are compared by name.} A record type \kw{sig\{\dots\}}
and a record element \kw{struct\{\dots\}} are compared by their occurrence
in the text: two occurrences of the same \kw{sig} are equal when the
variables they depend on have equal values, and two different occurrences
are different types even if their fields agree. This is the rule of
the proof checker \textsc{Half}~\cite{Ced} and of Gofer, and it avoids comparing large
records field by field. In practice a record comes with one named
constructor (Section~\ref{sec:segtypes}). Dependent pairs $(x:A)*B$ are
compared structurally.

\paragraph{The unfolding strategy.} \kw{opaque f} declares that, from this
point on, \kw{f} is not unfolded: it behaves as a constant, though its
body has not changed. \kw{transparent f} undoes this. The strategy is a
property of the checking run and not of the values: a value is computed
in the same way whatever the strategy, and only the comparison consults
it. In particular an \kw{opaque} declaration placed after a segment
applies to every instance of it. Making a definition opaque can only make
fewer terms convertible, so it never makes an incorrect proof correct.
The intended use is to lock a construction right after the lemmas that
characterise it, as one does in informal mathematics.

%======================================================================
\section{The module system}\label{sec:modules}

\subsection{Segments}

A segment is a named list of units:
\begin{verbatim}
  module Mon (A : U a) (e : A) (op : A -> A -> A) {
    twice (x : A) : A = op x x ;
    four (x : A) : A = twice (twice x) ;
  } ;
\end{verbatim}
The \emph{parameters} of a segment are its units without a body: the
telescope in the header and the parameters declared in the body. The
level variables and the constraints of its body belong to it as well. A
segment is checked once, where it is declared, with its parameters as
unknowns.

An \emph{instance} \kw{Mon[N, zero, plus]} gives values to the parameters,
in order. Each argument is checked against the type of its parameter, in
which the earlier arguments have been substituted; a level parameter
takes a level; a constraint of the segment must hold for the given
levels. Members are then accessed as \kw{Mon[N, zero, plus].four}, or
through an alias \kw{module NM = Mon[N, zero, plus] ;} as \kw{NM.four}, or
after \kw{open Mon[N, zero, plus] ;} as \kw{four}. Arguments may also be
given by name, \kw{Mon[A := N, \dots]}.

Segments nest. A segment declared inside another sees the parameters of
the enclosing one lexically, and needs no arguments for them:
\begin{verbatim}
  module Outer (A : U a) {
    module Inner { k (x : A) : A = x ; } ;
    kk (x : A) : A = Inner.k (Inner.k x) ;
  } ;
\end{verbatim}
\kw{import F ;} reads the file \kw{F.seg} and inserts it at that point as
a segment \kw{module F \{ \dots \} ;}, once. An import inside a segment
therefore makes the file a sub-segment, which sees the enclosing
parameters.

\subsection{Instances share}

An instance does not copy the body of the segment. It only records the
values of the parameters, in an environment; the members are evaluated
in that environment when they are used. A member therefore keeps a single
identity, its qualified name, and two occurrences of a member are
compared by this name and the values of the parameters it actually
depends on. In particular \kw{Mon[N, zero, plus].twice} reached through two
different aliases or two different \kw{open}s is one object, and the
equation
\begin{verbatim}
  share : Id (N -> N) (Mon[N, zero, plus].twice) (NM.twice) = refl NM.twice ;
\end{verbatim}
holds without unfolding \kw{twice}. When a large development instantiates
the same segment many times, as a formalisation of a model typically does
(a construction made over an abstract base, then applied to each
instance), this is what keeps checking fast: equal things are recognised
as equal by their names.

\subsection{Parameters stated once}\label{sec:params}

Because segments nest and share, assumptions of a whole development are
stated once, as parameters of one outer segment. For instance the model
formalised with \sys assumes function extensionality and univalence for
the metatheory; the main file reads
\begin{verbatim}
  module Model (fe : [b] Axioms[b].FunExt) (univ : [b] Axioms[b].Univalence) {
    import V ;
    import Presheaf ;
    ...
  } ;
\end{verbatim}
and every imported file sees \kw{fe} and \kw{univ}. Likewise the abstract
base of the model (a category with an interval and a presheaf of
cofibrations) is one segment whose parameters are the data of the base;
the generic constructions are imported inside it, and an instance of the
base (posets, or simplicial sets) is one instance of that segment, after
which every generic construction is a member of the instance. Excluded
middle, used at one place, is the parameter of a separate segment, so the
files that may use it are exactly those imported inside that segment.

\paragraph{Changing the assumptions.} Function extensionality follows from
univalence. Once this is proved, as a segment \kw{FunExtUA} producing, from
univalence at a level $b$, function extensionality at $b$, the whole
development is obtained under univalence alone by one instance:
\begin{verbatim}
  module ModelUA (univ : [b] Axioms[b].Univalence) {
    feU : [b] Axioms[b].FunExt = <b> FunExtUA[b].FE[univ @b].fe ;
    module M = Model[feU, univ] ;
  } ;
\end{verbatim}
Every result of the development is then available as
\kw{ModelUA[univ].M.\dots}. Since an instance only binds the parameters,
checking it checks the two arguments and nothing of the body: the
development is neither copied nor checked again.

\subsection{Segments as types}\label{sec:segtypes}

A segment is also a \emph{type}: the type of its instances, that is, the
dependent record of its parameters. Its defined members are available on
every element, as methods. With
\begin{verbatim}
  module Mon (l : Level) {
    A : U l ;  op : A -> A -> A ;
    sq (x : A) : A = op x x ;
    e : A ;  sqe : Id A (sq e) e ;
  } ;
\end{verbatim}
\begin{itemize}
\item \kw{Mon[a]} is a type (the level parameters must be given), and
  \kw{new Mon[a, N, plus, zero, refl zero]} is an element;
\item for \kw{M : Mon[a]}, \kw{M.op} is a field and \kw{M.sq} is the member
  evaluated in \kw{M}; \kw{open M} gives them short names;
\item a partial instance \kw{Mon[a, N]} or \kw{Mon[l := a, A := N]} is the
  type of the elements with those parameters fixed; for \kw{M : Mon[a, N]},
  \kw{M.A} computes to \kw{N};
\item fixing more parameters gives a subtype: \kw{Mon[a, N]} is a subtype
  of \kw{Mon[a]};
\item elements have $\eta$: \kw{M} is equal to
  \kw{new Mon[a, M.A, M.op, M.e, M.sqe]};
\item the universe of \kw{Mon[a]} is the maximum of the universes of its
  parameters.
\end{itemize}
A segment whose body has fields comes with a constructor \kw{mkMon},
defined right after it. Since elements are compared by occurrence, as
records are, using this one constructor everywhere makes elements built
from equal data equal. This gives records in the style of Arend's classes
and of \textsc{Half}'s theories, within the same mechanism as modules.

%======================================================================
\section{The checker}

The checker is written in Haskell, about 4500 lines, with no library
beyond the base. Its parts are: levels and the decision of constraints;
evaluation into values with closures; head normalisation under the
unfolding strategy; conversion; type checking of units; the filling of
holes; and the parser and scope resolution, which also computes, for each
definition, the variables it depends on.

The context after a unit does not depend on whether the unit checked:
definitions enter the context by their types and bodies, not by their
proofs of correctness. Hence the result of each declaration can be
computed separately, with a time budget (ten seconds per declaration;
an overrun counts as a failure), and one file can be checked while the
others are trusted. The checker also runs as a server that loads a
development once and, after an edit, checks only the files whose content
changed and those depending on them; the 205 files of the model check
from scratch in about 28 seconds.

\begin{remark}
What \sys does not have: no unification and no implicit arguments beyond
the matching of Section~\ref{sec:hidden}; no indexed inductive families;
no termination or positivity check; no pattern matching on \kw{refl}; no
proof irrelevance. These omissions are deliberate: each would make
checking less predictable, and the formalisation did not need them.
\end{remark}

%======================================================================
\section{Notation read before checking}\label{sec:notation}

The following conveniences are not part of the core language: they are
translated into core terms before checking, and the core never sees them.

\subsection{Omitted arguments}\label{sec:hidden}

A parameter written in braces, \kw{f \{x : A\} (y : B) : C}, is
\emph{hidden}: the type of \kw{f} is the ordinary
$(x:A)\to(y:B)\to C$, and a use \kw{f u} is read as \kw{f \_ u}, where
\kw{\_} is a hole. A hole may also be written explicitly anywhere, and a
hidden argument may be given explicitly as \kw{\{e\}}. Holes are filled
before the term is checked, by one-way \emph{matching}, never by
unification:
\begin{itemize}
\item the inferred type of each explicit argument is matched against the
  corresponding domain, with the holes as unknowns; two applications are
  matched argument by argument only when they have the same head;
  otherwise both are computed and matching continues on their head normal
  forms;
\item the expected type only fills holes left open by the explicit
  arguments;
\item two different solutions for one hole are an error, so the result
  does not depend on the order;
\item a solution is written back as a term in the current scope, using
  names only.
\end{itemize}
A hole that is not determined this way is an error that asks for an
explicit \kw{\{e\}}. This works well for lemmas about paths, such as
\kw{transH p (symH q)}, whose endpoints are determined by the explicit
arguments, and it is deliberately not clever.

\subsection{Equational chains} The notation
\begin{verbatim}
  calc A { a0 =< p1 > a1 =< p2 > a2 ... =< pn > an }
\end{verbatim}
proves $\Id\,A\,a_0\,a_n$. It stands for the right-nested composite of
the paths $p_i$; each $p_i$ is checked against $\Id\,A\,a_{i-1}\,a_i$, so
every point is written once and an error is reported at its step.



\subsection{Scoping}
\kw{open M ;} (or \kw{open M[args] ;}) makes the members of a segment
available by their short names. The resolver replaces every short name by
the qualified name it stands for, \kw{M[args].x}, so the core only sees
qualified names. For a variable \kw{c} of a segment type, \kw{open c} gives
short names to \kw{c.x}. \kw{import F ;} is expanded when the files are
read: the file \kw{F.seg} is inserted as \kw{module F \{ \dots \} ;}, once.
A \kw{where} block is a \kw{let}, and the constructor \kw{mkM} of a segment
type (Section~\ref{sec:segtypes}) is an ordinary definition added by the
resolver after the segment.


\begin{thebibliography}{9}
\bibitem{BC} M.~Bezem and T.~Coquand. Loop-checking and the uniform word
  problem for join-semilattices with an inflationary endomorphism.
  \emph{Theoretical Computer Science} 913, 2022.
\bibitem{BCDE} M.~Bezem, T.~Coquand, P.~Dybjer and M.~Escard\'o. Type theory
  with explicit universe polymorphism. arXiv:2212.03284, 2022 (revised and
  extended version, 2024).
\bibitem{Ced} J.~Cederquist. \emph{A Pointfree Approach to Constructive
  Analysis in Type Theory}. PhD thesis, Chalmers University of Technology and
  University of G\"oteborg, 1997. Section~7.1 describes the proof checker
  \textsc{Half}.
\bibitem{Wang} Q.~Wang. \emph{A Haskell Implementation for a Dependent Type
  Theory with Definitions}. Master's thesis, Chalmers University of Technology
  and University of Gothenburg, 2021.
\end{thebibliography}

\end{document}
